\scanpage{46}
la composante ${\mathfrak g}^{j}_{\deleted{\unsure{\kappa}}}$ ($1\leq j\leq h$) de $\mathfrak g$ suivant le caractère $\chi_j$ est
\[
\sum_{\substack{\alpha\in R\\o(\alpha)=j}}\mathfrak g^{\deleted{\unsure{j}}}_{\alpha}
+\sum_{\substack{\alpha\in R\\o(\alpha)=j-h}}\unsure{\widetilde{\mathfrak g}|\mathfrak g}_{\alpha}.
\]
En particulier $\mathfrak g^{1}$ est de rang $r+1$ ($r=$ rang).

Le centralisateur\inserted{$G^{\unsure{\kappa|K}}=N(T)^{\kappa}$} de $\kappa_{\unsure{h}}$ est un sous-groupe lisse de $N(T)$, image inverse du sous-groupe $\deleted{\unsure{W'|W}}$ de $W_{\kappa}$, \deleted{centralisateur} stabilisateur de $\kappa$ dans $W$. Kostant dit que le rang de $W_{\kappa}$ est égal à celui du centre de $\widetilde G$, et on peut espérer qu'il soit isomorphe à $\overline{P/R}$, $P=\widehat T=$ réseau des poids, $R=$ réseau des racines, \deleted{\unsure{ou}} ou à son dual \deleted{\unsure{de T}} à valeurs \unsure{ds|sur|dans} $\mathbb Q/\mathbb Z$.

La donnée d'un
\[
\kappa:\mu_h\unsure{-|\to}T
\]
de Kostant dans un tore maximal définit sur le système\inserted{de racines $R(T)$ de ce} dernier un système de racines \unsure{Q-simple|Qu-simple|Q-simples}\inserted{$\Sigma'$}, ou \unsure{Q-base|Qu-base} : l'ens des racines $\alpha$ telles que $\mathfrak g_{\alpha}\subset\mathfrak g^1$ ($\mathfrak g^1$ désigne la composante de $\mathfrak g$ suivant le car. $\chi_1$ de $\mu_h$). \inserted{En fait,} $\kappa$ est connu quand on connaît $\deleted{\unsure{Q}}\Sigma'\subset R$. \deleted{\unsure{Pb} des signes si $r>1$ d'où $h>2$ (donc la symétrie de $\mu_h$ n'est pas l'identité)} Il revient au $=$, dans le cas déployé, de dire que le \deleted{centralisateur} stabilisateur\inserted{$W_{\kappa}$} de $\kappa$ dans $W$ est \uline{égal} (à priori il est $\subset$) le stabilisateur de $\Sigma'$ dans $W$, i.e. l'ens des $w\in W$ tels que $w(\Sigma)\subset\Sigma'$, i.e. tels que $w$ transforme \uline{toutes} les racines de $\Sigma$ en
% archive sheet number: 45

